\documentclass[11pt]{article}

\usepackage[margin=1in]{geometry}
\usepackage[T1]{fontenc}
\usepackage[utf8]{inputenc}
\usepackage{lmodern}
\usepackage{microtype}
\usepackage{array}
\usepackage{booktabs}
\usepackage{tabularx}
\usepackage{longtable}
\usepackage{ragged2e}
\usepackage{enumitem}
\usepackage[hidelinks]{hyperref}
\usepackage{parskip}

\urlstyle{same}
\setlist[itemize]{leftmargin=1.3em}
\setlength{\emergencystretch}{2em}
\renewcommand{\arraystretch}{1.2}

\newcolumntype{Y}{>{\RaggedRight\arraybackslash}X}
\newcolumntype{P}[1]{>{\RaggedRight\arraybackslash}p{#1}}

\title{Mapping ISO/IEC/IEEE 42010:2022 to the SEI Views and Beyond Approach}
\author{}
\date{May 7, 2026}

\begin{document}
\maketitle

\begin{center}
\begin{minipage}{0.94\textwidth}
\small
\textbf{Purpose.} This handout provides a defensible crosswalk between ISO/IEC/IEEE 42010:2022 and the SEI \emph{Views and Beyond} (V\&B) documentation approach. It is written for teams that want to satisfy the requirements of an ISO 42010 architecture description while using SEI's practical view-based documentation method.\cite{iso42010,conceptualmodel,vandb}
\end{minipage}
\end{center}

\section*{Executive summary}
ISO/IEC/IEEE 42010:2022 defines requirements for the structure and expression of an architecture description, including architecture viewpoints and model kinds, but it does \emph{not} prescribe a single documentation method, notation set, or fixed list of viewpoint names.\cite{iso42010,conceptualmodel} The SEI V\&B approach provides a practical way to realize those requirements by selecting view types, choosing suitable styles and notations, and documenting rationale and cross-view relationships.\cite{iso42010,vandb}

\textbf{Important correction.} Labels such as \emph{functional/logical}, \emph{runtime}, \emph{deployment}, and \emph{development/build} are best treated as \textbf{common project-defined viewpoints in an ISO 42010-compliant architecture description}, not as viewpoint names defined normatively by ISO 42010 itself.\cite{iso42010,conceptualmodel}

\section*{Practical mapping}
\begingroup\footnotesize
\begin{longtable}{@{}P{0.20\textwidth} P{0.17\textwidth} P{0.28\textwidth} P{0.15\textwidth}@{}}
\toprule
\textbf{Common project-defined viewpoint in an ISO 42010 AD} & \textbf{SEI V\&B view type most often used} & \textbf{Typical stakeholder concerns addressed} & \textbf{Typical SEI styles or representations} \\
\midrule
\endfirsthead
\toprule
\textbf{Common project-defined viewpoint in an ISO 42010 AD} & \textbf{SEI V\&B view type most often used} & \textbf{Typical stakeholder concerns addressed} & \textbf{Typical SEI styles or representations} \\
\midrule
\endhead
Functional or logical & Module view type & Functional responsibility, decomposition, modifiability, reuse, separation of concerns, and build-time structure & Decomposition, Uses, Layered, Generalization \\
Runtime or interaction & Component-and-Connector (C\&C) view type & Runtime behavior, coordination, performance, reliability, availability, security, and data flow & Client-Server, Publish-Subscribe, Peer-to-Peer, Shared Data \\
Deployment or physical allocation & Allocation view type & Distribution, infrastructure placement, hardware constraints, resource utilization, cost, and operational environment & Deployment, Install \\
Development or work organization & Allocation view type & Team ownership, work partitioning, build responsibility, and project coordination & Work Assignment \\
\bottomrule
\end{longtable}
\endgroup

\section*{How the alignment works}
\begin{itemize}
  \item \textbf{ISO 42010 sets the obligation.} The standard requires an architecture description to identify stakeholders, concerns, viewpoints, model kinds, and related architectural concepts, while distinguishing the architecture itself from the architecture description that expresses it.\cite{iso42010,conceptualmodel}
  \item \textbf{V\&B supplies the documentation method.} V\&B organizes architecture documentation around selected views and supplements them with rationale, background information, and cross-view mappings.\cite{vandb}
  \item \textbf{Project viewpoints are concern-driven.} In an ISO 42010-compliant description, the architect defines viewpoints to frame stakeholder concerns. In practice, many software teams use project labels such as logical, runtime, or deployment as shorthand for those viewpoint definitions.\cite{iso42010,conceptualmodel}
  \item \textbf{SEI view types operationalize those viewpoints.} Module, C\&C, and Allocation views are often the most natural V\&B choices for expressing the models needed to satisfy those concern sets.\cite{vandb}
\end{itemize}

\section*{What to say precisely}
The following wording is accurate and review-friendly:

\begin{quote}
ISO/IEC/IEEE 42010:2022 specifies what an architecture viewpoint must accomplish within an architecture description, including its role in addressing stakeholder concerns through model kinds. The SEI \emph{Views and Beyond} approach provides a practical documentation method --- view types, styles, templates, and rationale guidance --- that can be used to realize those requirements in software architecture documentation.\cite{iso42010,conceptualmodel,vandb}
\end{quote}

\section*{What not to overclaim}
\begin{itemize}
  \item Do not state that ISO 42010 defines a mandatory canonical list of viewpoints named functional, runtime, deployment, or development. It does not.\cite{iso42010,conceptualmodel}
  \item Do not state that SEI styles are identical to ISO model kinds. A safer statement is that SEI styles and chosen notations can \emph{realize} the model kinds selected for a viewpoint.\cite{iso42010,conceptualmodel,vandb}
  \item Do not treat the crosswalk table as normative. It is a practical mapping used to build a compliant architecture description, not a standard-issued equivalence table.\cite{iso42010,conceptualmodel}
\end{itemize}

\section*{Recommended application workflow}
\begin{enumerate}[leftmargin=1.5em]
  \item Identify the architecture stakeholders and their concerns.\cite{iso42010,conceptualmodel}
  \item Define the project viewpoints needed to address those concerns and decide what model kinds each viewpoint requires.\cite{iso42010,conceptualmodel}
  \item Select the V\&B view type or types that best express those models.\cite{vandb}
  \item Choose the concrete style or notation for each view and produce the actual views.\cite{vandb}
  \item Document rationale, assumptions, mappings across views, and other ``beyond views'' material so the architecture description remains analyzable and usable.\cite{vandb}
\end{enumerate}

\section*{Bottom line}
Use ISO/IEC/IEEE 42010:2022 to define the \emph{requirements} for the architecture description and use SEI V\&B to define the \emph{documentation strategy}. That combination preserves standards alignment while giving the team a concrete and proven way to structure software architecture documentation.\cite{iso42010,conceptualmodel,vandb}

\begin{thebibliography}{9}

\bibitem{iso42010}
International Organization for Standardization.
\newblock \emph{ISO/IEC/IEEE 42010:2022 --- Software, systems and enterprise --- Architecture description}.
\newblock ISO, Edition 2, November 2022.
\newblock Official standard page: \url{https://www.iso.org/standard/74393.html}.
\newblock Accessed May 7, 2026.

\bibitem{conceptualmodel}
ISO/IEC/IEEE 42010 Website.
\newblock \emph{ISO/IEC/IEEE 42010: Conceptual Model}.
\newblock \url{https://www.iso-architecture.org/42010/cm/}.
\newblock Accessed May 7, 2026.

\bibitem{vandb}
Clements, P.; Bachmann, F.; Bass, L.; Garlan, D.; Ivers, J.; Little, R.; Merson, P.; Nord, R.; and Stafford, J.
\newblock \emph{Documenting Software Architectures: Views and Beyond}, 2nd ed.
\newblock Addison-Wesley, 2010. ISBN 978-0321552686.
\newblock Bibliographic details cited by Carnegie Mellon Software Engineering Institute in \emph{A Holistic View of Architecture Definition, Evolution, and Analysis} (CMU/SEI-2023-TR-004): \url{https://insights.sei.cmu.edu/documents/5720/2023_005_001_983542.pdf}.

\end{thebibliography}

\end{document}
